% Chapitre : sensibilité et incertitude à 2050 (axe 5).
% Sources principales : v2/docs/lab/reports/2026-11-14-axe5-sensibilite.md, v2/papers/axe5-sensibilite/main.tex,
% v2/sensibilite/ (code), v2/docs/lab/decisions.md (entrées « Axe 5 » du 2026-11-14 et « Axe 3, nuit »).
% Les valeurs anciennes qui circulent encore dans le dépôt sont relevées en annexe (ann:incoherences).
\begingroup\sloppy\setlength{\emergencystretch}{3em}% coupures des chemins longs, local à ce fichier

\section{Sensibilité et incertitude à 2050}\label{sec:sensibilite-ref}\label{sen:sec}

L'analyse de sensibilité (axe~5) répond à deux questions. Quels facteurs déplacent l'effectif privé en 2030, 2040
et 2050, et de combien ? Le choix du décideur pèse-t-il plus, dans la variance de cet effectif, que ce que l'on
ignore des données et des paramètres ? Elle sert aussi à mettre un prix sur l'axe~1 : il est inutile de sourcer
finement un paramètre qui ne déplace pas la sortie. Le code est dans \texttt{v2/sensibilite/} ; le rapport est
\texttt{v2/docs/lab/reports/2026-11-14-axe5-sensibilite.md} et le prétirage \texttt{v2/papers/axe5-sensibilite/main.tex}.
Toutes les valeurs de ce chapitre sont mesurées sur le jumeau (chapitre~\ref{jum:sec}), dans sa composition de
référence A0 et non avec l'estimateur servi C1, sauf mention contraire. Le facteur « filière » a été mesuré le 14 novembre avec
un fichier d'inscrits fourni à l'équipe sous accord de confidentialité ; depuis le 15, il lit SISE (données
ouvertes), et les parties~1 et~2 ont été relancées pour les dentistes et les pharmaciens, seules professions
concernées (section~\ref{sen:sise}).

\subsection{Relance du 15 novembre : vivier SISE pour les dentistes et les pharmaciens}\label{sen:sise}
Le 15 novembre 2026, le vivier de formation des dentistes et des pharmaciens est passé d'un fichier fourni sous
accord de confidentialité aux données ouvertes SISE (chapitre~\ref{sec:repro-ref}). SISE reproduit ce fichier à
l'identique pour les rentrées 2020 à 2024\mesure, mais n'a pas encore la rentrée 2025, où la 2e année de dentaire
passait de 1\,546 à 2\,083 inscrits\mesure. Les parties~1 et~2 ont été relancées pour ces deux professions
($n = 2\,048$, $r = 30$, même graine ; \texttt{v2/sensibilite/out/*\_sise.csv}, 526\,s\mesure).

\begin{center}\small
\begin{tabular}{@{}llrrl@{}}
\toprule
Profession & Famille & $S_T$(filière), ancien fichier & $S_T$(filière), SISE & Fourchette 2050 (p05--p95), SISE \\
\midrule
Dentistes & tout & 0,444\mesure & 0,020\mesure{} [0,019 ; 0,021] & 59\,932 -- 90\,964\mesure \\
Dentistes & incertitude & 0,547\mesure & 0,029\mesure{} [0,027 ; 0,031] & 60\,274 -- 84\,991\mesure \\
Pharmaciens & tout & 0,226\mesure & 0,045\mesure{} [0,042 ; 0,048] & 64\,163 -- 95\,043\mesure \\
Pharmaciens & incertitude & 0,320\mesure & 0,068\mesure{} [0,065 ; 0,072] & 64\,935 -- 88\,939\mesure \\
\bottomrule
\end{tabular}
\end{center}

\paragraph{Lecture.} Le poids du facteur « filière » venait presque entièrement d'une seule observation, la rentrée
2025, que la variante prolonge telle quelle jusqu'en 2050. Sans elle, la part privée et les quotas redeviennent les
premiers facteurs pour les dentistes ($S_T$ 0,377 et 0,325 en famille « tout »\mesure), et la fourchette 2050
des dentistes se resserre (p95 de 108\,797 à 90\,964\mesure). Aucune des deux versions n'est « la bonne » : la hausse
de 2025 est réelle, mais la prolonger 25 ans est une hypothèse. Le résultat à retenir est la fragilité de la
variante filière, pas l'une ou l'autre valeur. Les autres professions ne lisent pas ce vivier et ne changent pas.

\paragraph{Ce que ce chapitre ne donne pas.} Aucun chiffre 2050 n'est citable comme projection. Le registre l'a
conclu le 2026-11-14 après les protocoles de l'axe~3 : aucune bande n'est calibrée à son niveau nominal, et les
entrées du moteur sont biaisées vers le haut (\texttt{v2/docs/lab/decisions.md}, « Axe 3, nuit : \dots{} aucun
chiffre 2050 citable »). Aucun backtest n'existe au-delà de $h = 10$\mesure{}~ans (5\mesure{}~ans pour les
kinésithérapeutes et les infirmiers). Les niveaux et fourchettes ci-dessous décrivent la réponse du modèle à ses
facteurs ; ce ne sont pas des prévisions.

% ---------------------------------------------------------------------------------------------------
\subsection{Méthode}\label{sen:methode}

\paragraph{Sortie.} $Y_{p,t}$ est l'effectif privé national de la profession $p$, somme de
l'équation~\eqref{jum:eq:sortie} sur les départements, pour $t \in \{2030, 2040, 2050\}$. Le jumeau tourne au
département. L'origine est 2026, sauf pour les infirmiers (2023).

\paragraph{Deux parties.}
\begin{itemize}
  \item \textbf{Partie 1} (stock-flux de référence, privé seul). Sans régulation ni transition, aucun flux ne
        relie le privé et le public. Pour les infirmiers, la table par secteur n'a ni âge ni sexe : la partie~1
        prend les infirmiers libéraux (SNDS), soit \num{98854}\mesure{} en 2023 contre \num{277788}\mesure{}
        dans tout le privé, environ 36\derive{}\,\%.
  \item \textbf{Partie 2} (composition \texttt{LEVIERS}, leviers branchés). Pour les infirmiers, le périmètre
        devient le privé complet, sous une hypothèse d'indépendance secteur--âge--sexe (2013--2023) qui n'est pas
        variée. Le statu quo égale le tendanciel du jumeau à une personne près\mesure{}.
\end{itemize}

\paragraph{Facteurs de la partie 1} (tableau~\ref{sen:tab:facteurs}). Les niveaux discrets sont
équiprobables\suppose{}. Deux familles d'analyse : « tout » (leviers compris) classe ce qui déplace $Y$ ;
« incertitude » (leviers au tendanciel) donne la fourchette~(a). La variante d'entrées « tendance » (pente
prolongée) est exclue du plan : elle multiplie $Y_{2050}$ par 4 à 10\mesure{} avec une fenêtre de 2~ans
(sages-femmes : \num{183200} contre \num{19000}\mesure{}).

\begin{table}[h]
\centering\small
\begin{tabular}{p{4.0cm}llp{3.6cm}}
\toprule
Facteur & Famille & Plage & Statut de la plage \\
\midrule
fenêtre d'ajustement & estimation & $\{2, 3, 5, 8\}$ ans & choix \\
\texttt{shrink} & estimation & $\{20, 200, 2000, \text{national}\}$ & choix \\
entrées & estimation & \{plates, récentes 2 ans\} & choix \\
filière (CD, PH) & estimation & \{non, oui\} & choix ; inscrits DREES 2020--2025 (fichier sous accord de
  confidentialité) \\
part privée des salariés ambigus (sauf IDE) & données & \{min, centrale, max\} & bornes mesurées (axe~4) \\
fraction de vieillissement & paramètre & $[0{,}18 ; 0{,}22]$ & supposé\suppose{} \\
part en cumul emploi-retraite & paramètre & fiche $\times [0{,}5 ; 2]$ & sourcé (COR 2016) $\times$ supposé \\
report de la retraite & levier & $[0 ; 2]$ ans & levier \\
multiplicateur de cumul & levier & $[1 ; 2]$ & levier \\
quotas (après $L_p$) & levier & $[0{,}8 ; 1{,}2]$ & levier ; DD76 $\pm 20$\,\% \\
\bottomrule
\end{tabular}
\caption{Facteurs de la partie~1 (\texttt{v2/papers/axe5-sensibilite/main.tex}, tableau~1).}
\label{sen:tab:facteurs}
\end{table}

\paragraph{Plans d'expérience.}
\begin{itemize}
  \item \textbf{Morris} \citep{morris1991} : $r = 30$ trajectoires, $p = 4$ niveaux, pas $\Delta = 2/3$ de la
        plage. On rapporte $\mu^*$, moyenne des effets élémentaires absolus \citep{campolongo2007}, en \% de $Y$
        tendanciel.
  \item \textbf{Sobol} \citep{sobol2001} : matrices $A$, $B$ et $A_B^{(i)}$, $N = 2048$, soit $N(k+2)$ runs. Avec
        $Y$ centré et $V$ sa variance,
        \begin{equation}\label{sen:eq:sobol}
          S_i = \frac{\overline{f(B)\,\bigl[f(A_B^{(i)}) - f(A)\bigr]}}{V}, \qquad
          S_{T,i} = \frac{\overline{\bigl[f(A) - f(A_B^{(i)})\bigr]^2}}{2V},
        \end{equation}
        $S_i$ par l'estimateur de \citet{saltelli2010}, $S_{T,i}$ par celui de \citet{jansen1999}. Intervalles à
        95\,\% par 500 rééchantillonnages bootstrap des lignes, graine \num{20261114}.
  \item \textbf{Groupes} (partie~2) : les indices de groupe s'obtiennent en remplaçant d'un bloc les colonnes du
        groupe ; $N = 512$ par profession, soit \num{9216} à \num{11765}\mesure{} runs, 20 à 58\mesure{}~min par
        profession sur 5 processus.
\end{itemize}
La partie~1 a demandé \num{201276}\mesure{} runs en \SI{1187}{\second}\mesure{} sur une machine chargée, soit 2,7 à
4,2\mesure{}~ms par run au repos une fois les ajustements en cache.

\paragraph{Groupes de la partie 2.}
\begin{itemize}
  \item \textbf{Décision} : quotas et capacité de formation ; report de la retraite ; cumul emploi-retraite
        $\times [0{,}5 ; 2]$ ; départs avant 60~ans $\times [0{,}5 ; 1{,}5]$ ; abandon en formation infirmière
        $[0{,}056 ; 0{,}105]$\derive{} par an ; diplômés à l'étranger $\times [0{,}5 ; 2]$ ; part des zones
        régulées $[0 ; 1]$.
  \item \textbf{Connaissance} : les facteurs de la partie~1 hors leviers ; le flux annuel de diplômés à l'étranger,
        bas ou haut (dentistes : 665 ou \num{1300}\derive{} par an) ; la part contraignante
        $\theta \in [0{,}5 ; 1]$\suppose{}, la destination des entrants bloqués et la part dérogatoire.
  \item \textbf{Amplitude des leviers.} Référence « DD76 » : quotas et capacité $\times [0{,}8 ; 1{,}2]$, retraite
        0 à 2~ans, plage sourcée \citep{drees2021dd76} et retenue par l'utilisateur·rice
        (\texttt{decisions.md}, 2026-11-14, « l'amplitude de référence des leviers est celle de DD76 »). Variante
        « interface » : $\times [0{,}7 ; 1{,}5]$, retraite 0 à 3~ans, plages du prototype ; publiée comme
        sensibilité.
\end{itemize}

\paragraph{Fourchettes à 2050.} Deux composantes, jamais additionnées :
\begin{itemize}
  \item \textbf{(a) propagation} : quantiles 5 et 95\,\% de $Y$ sur les $2N$ tirages de la famille
        « incertitude » ;
  \item \textbf{(b) backtest extrapolé} : banc à origines glissantes (protocole N1 de l'axe~3,
        $h \in \{1, 3, 5, 8, 10\}$), erreur relative du total national ; pente $a$ de $|\text{erreur}| \propto h$
        par moindres carrés sans constante, puis $\pm 1{,}6\, a\, h$ : méthode A3 de l'observatoire (outil antérieur
        de l'équipe), censé couvrir environ 8 cas sur 10 ; cité, non reproduit.
\end{itemize}
Au département (partie~2), la largeur (a) est $(P_{95} - P_5)/Y^{\text{tend}}$ sur 500 tirages de la connaissance
seule, décisions au statu quo. Elle est comparée à un seuil dérivé : deux fois l'effet national à 2050 de quotas
$\times 1{,}2$, mesuré par profession (\texttt{decisions.md}, 2026-11-14, « fourchettes départementales et seuil
\enquote{assez sûr pour décider} »). Sous ce seuil, l'effet de quotas $\times 1{,}2$ dépasse la demi-bande de
connaissance (pour les kinésithérapeutes et les dentistes, ce sont pourtant les diplômés à l'étranger qui dominent
la variance due aux leviers).

\paragraph{Version du code.} Le jumeau a changé pendant l'analyse. Les deux parties ont été relancées sur un
instantané figé du jumeau (\texttt{git archive} du commit \texttt{b64ce56}) et sur les données du 2026-11-14,
08:08 ; empreintes SHA-256 tronquées \texttt{c784b1a8} (secteur), \texttt{43ad2c58} (infirmiers),
\texttt{4e5da69f} (zonage) (\texttt{decisions.md}, 2026-11-14, « relance sur le jumeau b64ce56 »). Les sorties
antérieures sont gardées dans \texttt{v2/sensibilite/out/avant\_b64ce56/}.

% ---------------------------------------------------------------------------------------------------
\subsection{Résultats}\label{sen:resultats}

\paragraph{Partie 1 : sources de la variance en 2050.} Indices totaux $S_T$ (tableau~\ref{sen:tab:st}). Trois sources
portent l'essentiel :
\begin{itemize}
  \item la part privée des salariés ambigus (sages-femmes 0,78\mesure{}, pharmaciens 0,48\mesure{}) ;
  \item le lien entre formation et entrées : filière 0,55\mesure{} (dentistes) et 0,32\mesure{} (pharmaciens) ;
        en levier, les quotas pèsent de 0,19 à 0,61\mesure{} selon la profession (famille « tout ») ;
  \item la fenêtre d'ajustement : 0,53\mesure{} pour les kinésithérapeutes et les infirmiers.
\end{itemize}
En 2030, la part privée domine pour les quatre professions où elle est variée (0,69 à 0,95\mesure{}) ; pour les
infirmiers, la fenêtre (0,65\mesure{}) et les entrées (0,52\mesure{}).

\begin{table}[h]
\centering\small
\begin{tabular}{lp{5.6cm}p{5.6cm}}
\toprule
Profession & Famille « incertitude » & Famille « tout » \\
\midrule
Sages-femmes & part privée 0,78 & part privée 0,45 ; quotas 0,39 \\
Dentistes & filière 0,55 ; part privée 0,26 ; fenêtre 0,13 ; entrées 0,10 & filière 0,44 ; part privée 0,21 ;
  quotas 0,19 \\
Pharmaciens & part privée 0,48 ; filière 0,32 ; fenêtre 0,15 & part privée 0,37 ; quotas 0,25 ; filière 0,23 \\
Kinésithérapeutes & fenêtre 0,53 ; part privée 0,18 ; entrées 0,17 ; \texttt{shrink} 0,17 ; fraction 0,12 &
  quotas 0,61 [0,58--0,64] ; fenêtre 0,21 \\
Infirmiers (libéraux) & fenêtre 0,53 ; entrées 0,49 ; fraction 0,20 & quotas 0,53 [0,50--0,57] ; entrées 0,27 ;
  fenêtre 0,25 ; retraite 0,03 \\
\bottomrule
\end{tabular}
\caption{Partie~1 : indices totaux $S_T$ en 2050 (mesurés\mesure{}). Intervalles bootstrap dans
\texttt{v2/docs/lab/reports/2026-11-14-axe5-sensibilite.md}.}
\label{sen:tab:st}
\end{table}

\paragraph{Interactions et contrôles.} Sur $Y$ brut, $S_1$ sortait de $[0 ; 1]$ ($-0{,}34$ à $1{,}78$\mesure{}) ; sur
$Y$ centré, il reste dans $[-0{,}001 ; 1{,}008]$\mesure{}. Les écarts $S_T - S_1$ montrent des interactions : pour
les infirmiers, fenêtre $S_1 = 0{,}28$ contre $S_T = 0{,}53$\mesure{}, entrées $0{,}24$ contre $0{,}49$\mesure{} ;
pour les kinésithérapeutes, fenêtre $0{,}35$ contre $0{,}53$\mesure{}. Morris donne le même ordre : $\mu^*$ de 15 à
19\mesure{}\,\% de $Y_{2050}$ pour les quotas, 24\mesure{}\,\% pour la filière dentaire, de 4 à 8\mesure{}\,\% pour
la fenêtre et la fraction de vieillissement.

\paragraph{Résultat nul : le cumul emploi-retraite.} $\mu^* \le 0{,}25$\mesure{}\,\% de $Y_{2050}$ pour la part
comme pour le multiplicateur, $S_T < 0{,}01$\mesure{}. La question est fermée pour l'axe~1, pour le total national et dans ce modèle. Le report de la
retraite a un $S_T \le 0{,}02$\mesure{}, sauf pour les infirmiers en 2030 (0,14\mesure{}).

\paragraph{Partie 1 : fourchettes nationales à 2050} (tableau~\ref{sen:tab:fourchettes}).

\begin{table}[h]
\centering\small
\setlength{\tabcolsep}{4pt}
\begin{tabular}{lrlll}
\toprule
Profession & $Y^{\text{tend}}_{2050}$ & (a) P5--P95 & (b) backtest extrapolé & Biais à $h = 5$ \\
\midrule
Sages-femmes & \num{18991} & \num{16001}--\num{21000} ($-16$ / $+11$\,\%) & $\pm 66{,}2$\,\% & $+7{,}9$\,\% \\
Dentistes & \num{76397} & \num{64816}--\num{102482} ($-15$ / $+34$\,\%) & $\pm 34{,}6$\,\% & $-4{,}6$\,\% \\
Pharmaciens & \num{72309} & \num{65561}--\num{96092} ($-9$ / $+33$\,\%) & $\pm 8{,}7$\,\% & $-0{,}2$\,\% \\
Kinésithérapeutes & \num{160783} & \num{146683}--\num{176233} ($-9$ / $+10$\,\%) & $\pm 16{,}4$\,\% ($h_{\max} = 5$)
  & $-1{,}7$\,\% \\
Infirmiers (libéraux) & \num{115323} & \num{95647}--\num{125972} ($-17$ / $+9$\,\%) & $\pm 49{,}3$\,\%
  ($h_{\max} = 5$) & $+5{,}8$\,\% \\
\bottomrule
\end{tabular}
\caption{Réponse du modèle à ses facteurs, pas une projection citable. Fourchettes nationales à 2050
(partie~1). $Y^{\text{tend}}$, (a) et le biais sont mesurés\mesure{} ; les
niveaux discrets sont équiprobables\suppose{}. (b) est dérivé\derive{} avec un facteur 1,6 cité, non reproduit.}
\label{sen:tab:fourchettes}
\end{table}

Pour les sages-femmes, l'erreur de backtest est un biais qui croît, pas une dispersion : $+1{,}5$\mesure{}\,\% à
$h = 1$, $+7{,}9$\mesure{}\,\% à $h = 5$, $+18{,}9$\mesure{}\,\% à $h = 10$ (une seule origine). La moitié haute de leur bande (b) est donc
peu informative. Les deux composantes dépassent la tolérance de $\pm 3$\suppose{}\,\% fixée par le protocole N2 pour la
reproduction de DD76.

\paragraph{Partie 2 : décision contre connaissance} (tableau~\ref{sen:tab:groupes}).

\begin{table}[h]
\centering\small
\begin{tabular}{llll}
\toprule
Profession & 2050, DD76 (IC 95\,\%) & 2050, interface & 2030, DD76 \\
\midrule
Infirmiers (privé) & 0,89 [0,79--1,01] / 0,07 [0,03--0,11] & 0,93 / 0,01 & 0,87 / 0,08 \\
Kinésithérapeutes & 0,81 [0,71--0,92] / 0,12 [0,07--0,17] & 0,83 / 0,05 & 0,53 / 0,41 \\
Sages-femmes & 0,59 [0,50--0,69] / 0,38 [0,31--0,46] & 0,76 / 0,16 & 0,06 / 0,94 \\
Dentistes & 0,54 [0,45--0,63] / 0,50 [0,42--0,57] & 0,66 / 0,35 & 0,16 / 0,90 \\
Pharmaciens & 0,38 [0,31--0,46] / 0,61 [0,52--0,71] & 0,65 / 0,31 & 0,04 / 0,99 \\
\bottomrule
\end{tabular}
\caption{Partie~2 : part de la variance de l'effectif privé national due à la décision et à la connaissance
(décision / connaissance, indices de groupe du premier ordre, mesurés\mesure{}). L'interaction
$1 - S_{\text{déc}} - S_{\text{conn}}$ reste comprise entre $-0{,}04$ et $0{,}12$\mesure{}.}
\label{sen:tab:groupes}
\end{table}

\begin{itemize}
  \item En 2030, la connaissance l'emporte (0,90 à 0,99\mesure{}), sauf pour les infirmiers, où les leviers à délai
        court (capacité IFSI en 3~ans, abandon) agissent déjà, et pour les kinésithérapeutes (intervalles proches).
  \item En 2050, avec l'amplitude DD76, la décision pèse plus que la connaissance pour les infirmiers, les
        kinésithérapeutes et les sages-femmes. Les dentistes sont à égalité (intervalles chevauchants). Pour les
        pharmaciens, la connaissance l'emporte (0,61 contre 0,38\mesure{}).
  \item Avec l'amplitude de l'interface, la décision l'emporte en 2050 pour les cinq professions. Le verdict dépend
        donc de l'amplitude jouable retenue.
  \item Parmi les leviers (DD76, $S_T$ 2050), les diplômés à l'étranger dominent pour les kinésithérapeutes
        (0,50\mesure{}) et les dentistes (0,40\mesure{}) ; ailleurs, ce sont les quotas (0,11 à 0,39\mesure{}) et la
        capacité de formation (0,05 à 0,36\mesure{}). Report de la retraite $\le 0{,}01$, cumul 0,00, part des zones
        $\le 0{,}01$, $\theta$, destination et dérogation $\le 0{,}01$\mesure{} : ces leviers ne déplacent pas le total
        national.
\end{itemize}

\paragraph{Partie 2 : fourchettes départementales.} Effet national de quotas $\times 1{,}2$ en 2050 : sages-femmes
$+12{,}3$, dentistes $+11{,}9$, pharmaciens $+11{,}3$, kinésithérapeutes $+11{,}9$, infirmiers (privé)
$+13{,}7$\mesure{}\,\%. Le seuil vaut donc 22,6 à 27,4\derive{}\,\%. Le tableau~\ref{sen:tab:dep} donne la largeur
(a) départementale et le nombre de départements sous le seuil.

\begin{table}[h]
\centering\small
\begin{tabular}{lccc}
\toprule
Profession & Largeur (a) médiane [P10--P90] & (b) médiane & Départements sous le seuil \\
\midrule
Sages-femmes & 56\,\% [31--113] & $\pm 80$\,\% & 2 / 101 \\
Dentistes & 69\,\% [50--116] & $\pm 56$\,\% & 0 / 101 \\
Pharmaciens & 55\,\% [40--97] & $\pm 32$\,\% & 0 / 101 \\
Kinésithérapeutes & 34\,\% [19--54] & $\pm 46$\,\% & 22 / 101 \\
Infirmiers (privé) & 26\,\% [17--47] & $\pm 37$\,\% & 54 / 100 \\
\bottomrule
\end{tabular}
\caption{Partie~2 : largeur relative de la bande de connaissance à 2050 par département (mesurée\mesure{}), composante
backtest (dérivée\derive{}) et nombre de départements où la largeur (a) est inférieure au seuil dérivé.}
\label{sen:tab:dep}
\end{table}

Le seuil porte sur la composante (a) seulement : la composante (b) n'est pas observée au-delà de $h = 10$ (5 pour
les kinésithérapeutes et les infirmiers) ; en 2050, elle est extrapolée. Un département « sous le seuil » est un
département où l'effet de quotas $\times 1{,}2$ se distingue de la bande de connaissance ; ce
n'est pas un département où le niveau 2050 est connu.

\paragraph{Convergence.} Partie~1 ($N = 2048$) : la largeur maximale de l'intervalle de $S_T$ (facteurs
$S_T \ge 0{,}05$ en 2050) passe de 0,19--0,41 à $N = 128$ à 0,05--0,09\mesure{} à $N = 2048$ ; l'écart à la valeur
finale est $\le 0{,}021$\mesure{} dès $N = 1024$. Partie~2 ($N = 512$) : largeur maximale 0,15--0,19\mesure{}, sur
72 facteurs de $S_T \ge 0{,}05$ ; écart entre $N = 256$ et $N = 512$ $\le 0{,}06$\mesure{}, sauf pour les pharmaciens
(0,085\mesure{}). $N = 512$ suffit pour les parts de groupe, mais ni pour trancher le cas des dentistes, ni pour
départager deux leviers à $\pm 0{,}08$ près.

\paragraph{Où mettre l'effort.} Le gain attendu d'une étude est $S_T \times$ la fraction réductible (1 pour un choix
tranché par le banc, 0,75 pour une plage divisée par deux, 0,5 pour un encadrement resserré ; fractions
supposées\suppose{}). Les gains les plus grands portent sur la filière (dentistes 0,29, pharmaciens
0,17\derive{}) et sur la part privée des salariés ambigus (sages-femmes 0,15, pharmaciens 0,13\derive{}). Cumul,
$\theta$, destination et dérogation ont un gain $\approx 0$ : il est inutile de les approfondir pour le total national
(\texttt{v2/papers/axe5-sensibilite/main.tex}, tableau de priorisation).

% ---------------------------------------------------------------------------------------------------
\subsection{Limites}\label{sen:limites}

\begin{itemize}
  \item \textbf{Les plages font les indices.} L'équiprobabilité des niveaux discrets (une fenêtre de 2~ans pèse
        autant qu'une de 5) est supposée ; resserrer la fenêtre à $\{5, 8\}$ ferait baisser son indice. Les parts
        de groupe dépendent de l'amplitude jouable : le passage de DD76 à l'interface change le verdict des
        pharmaciens.
  \item \textbf{Double compte d'une même politique.} Quotas et capacité de formation sont deux mécanismes d'une même
        politique, comptés tous deux dans le groupe « décision ».
  \item \textbf{Infirmiers.} Partie~1 : libéraux seuls, environ 36\derive{}\,\% du privé. Partie~2 : privé complet
        sous une hypothèse d'indépendance secteur--âge--sexe non variée ; la part « connaissance » des infirmiers est
        donc sous-estimée.
  \item \textbf{Facteurs sans prise.} Dans la partie~1, la formation, la régulation (zones communales, non
        activables au département) et une vingtaine de fiches de l'axe~1 n'ont pas de point d'entrée : leur
        sensibilité est indécidable, pas nulle. Dans la partie~2, la régulation, $\theta$ et la destination n'agissent
        que sur la répartition entre communes, que l'analyse ne mesure pas : un $S_T$ nul ne dit rien de l'effet
        local.
  \item \textbf{Échelle.} Les indices portent sur le total national. Au département, le classement des facteurs
        peut différer (\texttt{shrink} y pèserait plus) ; il n'est pas mesuré.
  \item \textbf{Extrapolation.} La composante (b) prolonge linéairement des erreurs observées jusqu'à $h = 10$ (5 pour
        les kinésithérapeutes et les infirmiers) vers $h = 24$--27, avec un facteur 1,6 cité et non reproduit.
  \item \textbf{Modèle unique.} Les indices sont ceux du stock-flux v1 réécrit. Une méthode de l'axe~2 (diplômés
        $\to$ entrées, régulation communale) peut changer le classement.
  \item \textbf{Code en mouvement.} Avant la relance sur \texttt{b64ce56}, sages-femmes et dentistes penchaient du
        côté de la connaissance (0,41 contre 0,57--0,59\mesure{}) ; le verdict dépend donc de la version du jumeau.
        L'analyse n'a pas été relancée sur l'état du jumeau validé le 2026-11-15 (chapitre~\ref{jum:sec}). Deux
        défauts de l'analyse elle-même ont été corrigés avant la version publiée : l'année 2012 des infirmiers, qui
        ne contenait que les libéraux et produisait un saut fictif de $+190$\mesure{}\,\% ; et l'abandon de formation
        forcé à 0 au statu quo pour les sages-femmes et les kinésithérapeutes.
  \item \textbf{Pas de chiffre citable à 2050.} Les fourchettes ci-dessus ne sont pas des intervalles calibrés. Le
        registre conclut qu'aucun chiffre 2050 n'est citable et qu'aucune bande n'est publiable sous son niveau
        nominal (\texttt{decisions.md}, 2026-11-14, « Axe 3, nuit »).
\end{itemize}

\paragraph{Questions ouvertes chiffrées} (prix supposés\suppose{}, \texttt{v2/docs/lab/reports/2026-11-14-axe5-sensibilite.md}) :
analyse au département, environ 20~min de machine ; reproduire le facteur 1,6, environ 20~min ; plages pondérées par
l'erreur du banc N1, environ 30~min ; varier l'hypothèse d'indépendance des infirmiers, environ 30~min et 1~h de
machine ; Sobol par département, environ 2~h de machine ; doubler $N$ de la partie~2, environ 1~h 30 de machine.
\endgroup
